\section{Discussion}
\label{sec:discussion}

\subsection{Legal Theory: Why Isoperimetric Naturalness Matters}

GeoSection's central legal claim is that the natural ratio is determined by
geography, not by partisan composition.
The isoperimetric normalisation makes this claim auditable: any party can
independently compute the normalised edge-cut for every feasible ratio using
only TIGER boundary data and census populations.
If the algorithm selects $6:8$ for North Carolina, the justification is
publicly reproducible: $\mathrm{EC}(1:13)/\sqrt{1} = 98$\,km $>$ $\mathrm{EC}(6:8)/\sqrt{6} = 92.6$\,km.

This is a procedural rather than a partisan claim, which is important for
state-court review after \emph{Rucho v.\ Common Cause} (2019).
State courts applying their own anti-gerrymandering provisions have shown
receptiveness to procedural arguments: \emph{League of Women Voters v.\
Pennsylvania} (2018) and \emph{Harper v.\ Hall} (2022) both emphasised
the procedural neutrality of the algorithm used to construct the challenged
maps \citep{lwv2018pa,harper2022nc}.

\paragraph{Why the caterpillar is problematic.}
The caterpillar pathology is legally vulnerable for two reasons.

First, it systematically concentrates Democratic voters into a small number
of highly efficient districts while spreading Republicans across many
competitive ones.
This produces a proportionality gap mechanically, without any partisan input.
A court observer who does not understand the algorithm would see the same
pattern as in a deliberate partisan gerrymander: safe Democratic urban seats
ringing a sea of Republican-leaning suburban and rural districts.

Second, the caterpillar cannot be justified on geographic grounds.
If asked ``why did you split North Carolina $1:13$ instead of $6:8$?'',
the answer under standard bisection is simply ``because $1:13$ has a smaller
edge-cut.''
Under GeoSection, the answer is ``because the $6:8$ east/west split has a
smaller \emph{isoperimetrically normalised} edge-cut, confirming that it is
the natural geographic division.''
The second answer is far more defensible: it articulates a principled
geometric criterion that is neutral with respect to partisanship.

\paragraph{States where the peel is confirmed genuine.}
GeoSection does not eliminate all asymmetric splits.
For states where the urban core is genuinely geographically isolated, the
normalisation \emph{confirms} the peel rather than overriding it.
Philadelphia ($k=17$, PA): $\mathrm{EC}(1:16)/\sqrt{1} = 39$\,km $<$ any
normalised equal-split comparison.
The algorithm's selection of $1:16$ carries a stronger legal argument:
``we tried the equal split and it costs more per unit of geographic territory;
Philadelphia is genuinely compact and geographically isolated.''

This self-certifying property is important: GeoSection's answer is not ``we
imposed equal splits on principle'' but ``geography chose the ratio.''
For states where geography genuinely warrants a peel, the peel is confirmed.
For states where it does not (NC), a genuine bisection is selected.

\subsection{Separating Legal and Partisan Claims}

GeoSection makes two distinct and separable claims that should not be conflated.

\textbf{(a) Legal/geometric claim.}
The natural ratio is determined by geography alone: any party can independently
compute $\mathrm{EC}(i:k-i)/\sqrt{\min(i,k-i)}$ from public TIGER and census data.
This auditability is a legal merit regardless of whether the ratio selection changes
any partisan outcome.
GeoSection's selection of $6:8$ for North Carolina is defensible on geometric grounds
independent of the 5\,D/9\,R seat count.

\textbf{(b) Partisan claim.}
GeoSection changes the bisection tree structure.
In 6 of 44 states, this changes the seat count by $+1\,D$; in 1 state (Pennsylvania),
it changes by $-1\,D$.
In 38 states the seat count is unchanged.
The partisan claim is weak: GeoSection is not a proportionality fix.
Its primary legal merit is procedural (claim (a)), not partisan.

\paragraph{Pennsylvania: why MEC outperforms GeoSection.}
Pennsylvania is the one state where the MEC baseline produces a better Democratic
outcome than GeoSection.
GeoSection confirms the $1:16$ Philadelphia peel: Philadelphia's compact
130\,km boundary wins the isoperimetric comparison against any near-equal split.
Under MEC, the standard $8:9$ first-level bisection places several southeast
Pennsylvania suburban tracts (Delaware County, Chester County outer ring) in
the western bisection half.
These tracts lean Democratic ($\sim 54$--58\,\% D) and, when grouped with
western Pennsylvania's Pittsburgh-area districts, create one additional
lean-Democratic district in southwestern Pennsylvania.
Under GeoSection's $1:16$ confirmed peel, the Philadelphia subgraph (left half,
1 district) contains only Philadelphia proper; the Delaware County and Chester
County suburban tracts are placed in the 16-district right half where, after
recursive bisection, they are split between Republican-leaning rural districts
and a single Philadelphia suburb district that leans Republican at the margin
($\sim 51$\,\% R).
The net effect: the MEC suburban tract placement produces one additional
D-leaning district that GeoSection's confirmed peel does not.

This finding is not evidence of algorithmic bias.
It reflects the \emph{geographic sensitivity} of the suburban ring:
Chester and Delaware County's Democratic lean is marginal enough that
their placement within the bisection tree determines whether they tip
a district toward D or R.
The confirmed caterpillar does not create a better Republican outcome by design;
it creates a different geographic grouping that, for Pennsylvania's specific
partisan geography, happens to favour Republicans by one seat.

\paragraph{Confirmed caterpillar states (IL, PA): proportionality comparison.}
Illinois and Pennsylvania both receive GeoSection natural ratios of $1:16$,
confirmed by the isoperimetric normalisation.
The relevant question is whether this confirmed asymmetric structure produces
different proportionality outcomes than states where GeoSection fires the
normalisation shift.

\begin{table}[h]
\centering
\caption{Proportionality gap comparison: confirmed caterpillar vs.\ normalisation-shifted.}
\label{tab:caterpillar-comparison}
\begin{tabular}{llrrrrr}
\toprule
State & Category & $k$ & D vote & GEO D & MEC D & Gap (pp) \\
\midrule
\multicolumn{7}{l}{\textit{Confirmed caterpillar ($1:k-1$ wins normalisation)}} \\
IL & Confirmed & 17 & 58.7\,\% & 12 & 11 & $-5.5$ \\
PA & Confirmed & 17 & 50.6\,\% & 7  & 8  & $-9.1$ \\
\midrule
\multicolumn{7}{l}{\textit{Normalisation-shifted (equal split wins)}} \\
NC & Shifted & 14 & 49.3\,\% & 5  & 5  & $-13.6$ \\
WI & Shifted & 8  & 50.3\,\% & 3  & 2  & $-12.8$ \\
\bottomrule
\end{tabular}
\end{table}

The confirmed caterpillar states (IL, PA) do not show systematically worse
proportionality than the normalisation-shifted states (NC, WI).
Illinois, with the most extreme caterpillar ($1:16$), shows the smallest
gap ($-5.5$\,pp) because Chicago is genuinely one of the most Democratic cities
in the US and its isolation in a single district does not create excessive
ring-district waste.
Pennsylvania's $-9.1$\,pp gap reflects suburban Philadelphia's swing status,
not caterpillar structure per se.

The key normative implication: the legal distinction between a
\emph{confirmed} peel and a \emph{prevented} caterpillar does not produce
systematically different partisan outcomes.
GeoSection's isoperimetric confirmation changes the legal \emph{justification}
for a peel (``geography chose this ratio'') but not the partisan geography
that determines seat outcomes.

\subsection{Relationship to AreaSection (B.9)}

GeoSection and AreaSection (B.9) address the same caterpillar pathology but
through different mechanisms.

GeoSection says: ``among all feasible ratios, select the one whose boundary
is most compact per unit of geographic coverage.''
The constraint is on the \emph{selection criterion}: normalised edge-cut.

AreaSection says: ``require each bisection half to contain approximately
50\,\% of the state's land area.''
The constraint is on the \emph{partition itself}: dual population+area balance.

Both algorithms prevent the caterpillar pathology; they differ in how far
they go.
For North Carolina, both algorithms select $6:8$ as the natural ratio
(AreaSection via area constraint, GeoSection via normalised edge-cut).
For Illinois, GeoSection still selects $1:16$ (Chicago's compact boundary
wins the isoperimetric comparison), while AreaSection selects $8:9$
(Chicago occupies $<5$\,\% of Illinois's land, making area balance infeasible
at $1:16$).

Table~\ref{tab:geosection-vs-areasection} summarises the key differences.

\begin{table}[h]
\centering
\caption{GeoSection (B.8) vs.\ AreaSection (B.9) comparison.}
\label{tab:geosection-vs-areasection}
\begin{tabular}{lll}
\toprule
& GeoSection & AreaSection \\
\midrule
METIS mode & ncon=1 (population only) & ncon=2 (pop + area) \\
Constraint & Normalised edge-cut & Area balance ($\pm 10$\,\%) \\
IL first-level ratio & 1:16 & 8:9 \\
NC first-level ratio & 6:8 & 6:8 \\
WI first-level ratio & 1:7 & 3:5 (area forces higher) \\
Caterpillar prevention & Partial (density-dependent) & Stronger \\
Legal rationale & ``Compact per unit area'' & ``Geographic balance'' \\
Implementation & Simpler (ncon=1) & More complex (ncon=2) \\
\bottomrule
\end{tabular}
\end{table}

The two algorithms are complementary: GeoSection is simpler and faster;
AreaSection provides a stronger anti-caterpillar guarantee for the largest
states.
An analyst could reasonably choose GeoSection for states where the caterpillar
is not the primary concern, and AreaSection for states where area balance
is independently desirable (large, geographically diverse states).

\subsection{Why Partisan Outcomes Are Stable}
\label{sec:stability}

The head-to-head comparison in Section~\ref{sec:comparison} shows that
38 of 44 states produce identical seat counts under GeoSection and the
MEC baseline.
This stability is not a coincidence.

In states with strong geographic partisan sorting (urban Democrats, rural
Republicans), the \emph{number} of Democratic congressional seats is
determined by how many contiguous geographic regions are sufficiently
Democratic to elect a Democrat.
GeoSection changes the structure of the bisection tree --- which tracts
are in the same district --- but not which geographic regions are
sufficiently Democratic.

The B.7 seed-sweep convergence result establishes this rigorously for
the MEC case: for states where the MEC is unique, the partisan outcome
is determined by the graph topology.
GeoSection provides corroborating evidence: even when the bisection ratio
changes ($1:13$ to $6:8$ for NC), the partisan outcome does not change.
The five North Carolina districts with Democratic majorities are determined
by the locations of Charlotte, Raleigh, and the Research Triangle ---
not by whether they are all in the same bisection half.

The six states where GeoSection and MEC differ (CO, IL, MN, OK, WA, WI)
are cases where the specific placement of suburban ring tracts depends
sensitively on which bisection tree structure is used.
In Wisconsin, the MEC standard bisection places two suburban Milwaukee
tracts in the same district as Milwaukee's urban core, creating a 2\,D
outcome; GeoSection's $1:7$ split keeps Milwaukee as a compact single
district and distributes the suburban ring tracts differently, creating
one additional Democratic district in the Madison area.

\subsection{Limitations}
\label{sec:limitations}

\paragraph{Approximation of the natural ratio.}
The ratio scan approximates the true minimum edge-cut at each ratio using
$N=50$ seeds.
The natural ratio selected is the ratio with the minimum normalised edge-cut
\emph{among observed seeds}, which may differ from the true minimum.
For the seed-sweep states identified in B.7, 50 seeds is sufficient for
convergence; for large states (CA, TX, FL) with complex geography, higher
seed counts may be needed.
We mark these states as having uncertain natural ratios in Table~\ref{tab:natural-ratios}.

\paragraph{Isoperimetric approximation.}
Lemma~\ref{lem:isoperimetric} assumes approximately uniform density within
the state.
For states with extreme density concentration (Manhattan vs.\ upstate New York),
the approximation may not hold: the true isoperimetric cost of separating
Manhattan from the rest of New York depends on Manhattan's specific geographic
shape, not just the fraction of population it contains.
In practice, the normalisation works well because METIS reports the actual
edge-cut (not the approximated minimum), so the normalised comparison is
comparing real boundary lengths.

\paragraph{Phase 2 not evaluated.}
The directional penalty ($\lambda > 0$) is implemented but not included in
the 50-state sweep.
Its effect on boundary straightness and partisan outcomes remains to be
characterised.
Preliminary single-state tests suggest $\lambda = 0.1$--$0.3$ produces
noticeably straighter boundaries without changing seat counts.

\paragraph{Population balance in deep recursive levels.}
The per-node ufactor formula ($1 + \delta/k$) sets very tight tolerances
at the root level ($k=17$ for Illinois: $1.000294$) and relaxed tolerances
at leaf nodes ($k=2$: $1.0025$).
For the largest states (CA: $k=52$, TX: $k=38$), root-level ufactor is
extremely tight ($\sim 0.0001$\,\% tolerance).
METIS occasionally fails to meet this constraint, requiring the post-hoc
boundary-swap rebalancing.
In the 50-state sweep, all states pass the 0.5\,\% balance check at the
level used for the natural ratio; deeper recursive levels use relaxed
tolerances.

\paragraph{Single census year.}
All results use 2020 Census data.
The natural ratio structure may differ for 2000 and 2010 Census data
due to population shifts (urban growth, suburban expansion).
Cross-census stability of the natural ratio is deferred to future work.
